\section*{Aufgabe 1}

\subsection*{DPLL)}


$\varphi = (\neg X_0 \vee X_1 \vee X_2 ) \wedge
(X_0 \vee X_1 \vee X_3 )
\wedge
(X_0 \vee X_1 \vee \neg X_3 )
\wedge
(X_0 \vee \neg X_1 \vee X_3 )
\wedge
(X_0 \vee \neg X_1 \vee \neg X_3 )
\wedge
(\neg X_2 \vee \neg X_2 \vee X_3 )
\wedge
(\neg X_0 \vee \neg X_1 \vee X_2 )
\wedge
( \neg X_0 \vee X_1 \vee \neg X_2 )
$\\

$M= \{ \{ \neg X_0 , X_1 , X_2 \} 
\{ X_0 , X_1 , X_3 \}
\{ X_0 , X_1 , \neg X_3 \}
\{ X_0 , \neg X_1 , X_3 \}$

$
\{ X_0 , \neg X_1 , \neg X_3 \}
\{ \neg X_2 , \neg X_2 , X_3 \}
\{ \neg X_0 , \neg X_1 , X_2 \}
\{ \neg X_0 , X_1 , \neg X_2 \} \} $\\

Für Rekursion füge ich nach einer Schrittnummer ein $.$ hinzu und beginne wieder von 1 zu zählen falls ein Rückschritt notwendig sein könnte.
\subsubsection*{Schrit 1}
$M= \{ \{ \neg X_0 , X_1 , X_2 \} 
\{ X_0 , X_1 , X_3 \}
\{ X_0 , X_1 , \neg X_3 \}
\{ X_0 , \neg X_1 , X_3 \}$

$
\{ X_0 , \neg X_1 , \neg X_3 \}
\{ \neg X_2 , \neg X_2 , X_3 \}
\{ \neg X_0 , \neg X_1 , X_2 \}
\{ \neg X_0 , X_1 , \neg X_2 \} \} $\\
Es gilt $\emptyset \not \in M$ und $  \emptyset \neq M$ und $\{ L \} \not \in M$ und $ L  \not \in \cup M$ daher 5. Setze $X_0 = 1$ und $DPLL(M\vert_{X_0})$ 

\subsubsection*{Schrit 1.1}
$M= \{ \{ X_1 , X_2 \} 
\{ \neg X_2 , \neg X_2 , X_3 \}
\{ \neg X_1 , X_2 \}
\{ X_1 , \neg X_2 \} \} $\\
Es gilt $\emptyset \not \in M$ und $  \emptyset \neq M$ und $\{ L \} \not \in M$ und $ L  \not \in \cup M$ daher 5. Setze $X_1 = 1$ und $DPLL(M\vert_{X_1})$ 

\subsubsection*{Schrit 1.1.1}
$M= \{ 
\{ \neg X_2 , \neg X_2 , X_3 \}
\{  X_2 \}
 \} $\\
Es gilt $\emptyset \not \in M$ und $  \emptyset \neq M$ und $\{ L \}  \in M$ mit $L= X_2$ $DPLL(M\vert_{X_2})$ 

\subsubsection*{Schrit 1.1.2}
$M= \{ 
\{  X_3 \}
 \} $\\
Es gilt $\emptyset \not \in M$ und $  \emptyset \neq M$ und $\{ L \}  \in M$ mit $L= X_3$ $DPLL(M\vert_{X_3})$ 

\subsubsection*{Schrit 1.1.3}
$M= \{ 
 \} $\\
Es gilt $\emptyset \not \in M$ und $  \emptyset = M$ daher nach 2. $return \emptyset$\\

Als Ergebnis gibt es die Belegung $\{ X_0 \mapsto 1 , X_1 \mapsto 1 , X_2 \mapsto 1 , X_3 \mapsto 1  \}$

\subsection*{Sat}

 $\varphi_0 = \bigwedge (X_0 \nleftrightarrow X_1 , X_1 \nleftrightarrow X_2 , X_2 \nleftrightarrow X_0 )$\\


 $\varphi_1 = \bigwedge (X_0 \nleftrightarrow X_1 , X_1 \nleftrightarrow X_2 , X_2 \nleftrightarrow X_3 , X_3 \nleftrightarrow X_0 )$\\
 
 Baum ist auf der nächsten Seite, erklärung der Zeichen:\\
$\nleftrightarrow = ne$\\
$b = \bot$  und $t = \top$\\
$a = \wedge$
 
 
 
\includepdf{part/Serie5A1Zeich.pdf}
